\documentclass[11pt,a4paper]{article}
\usepackage[utf8]{inputenc}
\usepackage[T1]{fontenc}
\usepackage[french]{babel}
\usepackage[a4paper,top=1.75cm,bottom=1.9cm,left=1.85cm,right=1.85cm]{geometry}
\usepackage{amsmath,amssymb}
\usepackage{array}
\usepackage{booktabs}
\usepackage{enumitem}
\usepackage{eurosym}
\usepackage{fancyhdr}
\usepackage{hyperref}
\usepackage{lmodern}
\usepackage{inconsolata}
\usepackage{listings}
\usepackage{microtype}
\usepackage{tabularx}
\usepackage[most]{tcolorbox}
\usepackage{titlesec}
\usepackage{xcolor}

\renewcommand{\familydefault}{\sfdefault}

\definecolor{ink}{HTML}{202124}
\definecolor{muted}{HTML}{5F6368}
\definecolor{line}{HTML}{DADCE0}
\definecolor{paper}{HTML}{F8F9FA}
\definecolor{panel}{HTML}{FFFFFF}
\definecolor{accent}{HTML}{8A1538}
\definecolor{accentsoft}{HTML}{F8EEF2}
\definecolor{green}{HTML}{0B6B57}
\definecolor{greensoft}{HTML}{E6F4EF}
\definecolor{amber}{HTML}{9A5B00}
\definecolor{ambersoft}{HTML}{FFF7E6}
\definecolor{blue}{HTML}{245B8E}
\definecolor{bluesoft}{HTML}{EAF2FA}
\definecolor{rose}{HTML}{A84D77}
\definecolor{rosesoft}{HTML}{FAF0F5}
\definecolor{codeback}{HTML}{F8F9FA}
\definecolor{codemark}{HTML}{DADCE0}

\hypersetup{
  colorlinks=true,
  linkcolor=accent,
  urlcolor=green,
  pdfborder={0 0 0}
}

\setlength{\parindent}{0pt}
\setlength{\parskip}{5pt}
\setlist[itemize]{leftmargin=1.35em,itemsep=2pt,topsep=2pt}
\setlist[enumerate]{leftmargin=1.55em,itemsep=3pt,topsep=3pt}
\setlist[enumerate,1]{label=\textbf{\alph*.}}

\pagestyle{fancy}
\fancyhf{}
\fancyhead[L]{\sffamily\footnotesize Python pour économistes}
\fancyhead[R]{\sffamily\footnotesize Université de Lille -- L2 Économie}
\fancyfoot[C]{\sffamily\footnotesize\color{muted}\thepage}
\renewcommand{\headrulewidth}{0.2pt}
\renewcommand{\headrule}{\hbox to\headwidth{\color{line}\leaders\hrule height \headrulewidth\hfill}}
\setlength{\headheight}{22pt}

\titleformat{\section}
  {\sffamily\Large\bfseries\color{ink}}
  {}{0pt}{}
  [\vspace{-0.25em}{\color{accent}\titlerule[0.7pt]}]
\titleformat{\subsection}
  {\sffamily\large\bfseries\color{ink}}
  {}{0pt}{}

\lstdefinelanguage{terminal}{
  morekeywords={ls,cd,pwd,mkdir,python,python3,code,dir,type,echo},
  sensitive=true
}

\lstset{
  basicstyle=\ttfamily\footnotesize,
  backgroundcolor=\color{codeback},
  frame=single,
  framesep=5pt,
  framerule=0.25pt,
  rulecolor=\color{codemark},
  language=Python,
  keywordstyle=\color{accent}\bfseries,
  commentstyle=\color{muted},
  stringstyle=\color{green},
  numberstyle=\scriptsize\color{muted},
  numbers=left,
  numbersep=8pt,
  showstringspaces=false,
  breaklines=true,
  aboveskip=0.7em,
  belowskip=0.7em,
  xleftmargin=1.1em,
  framexleftmargin=1em,
  columns=fullflexible,
  keepspaces=true,
  upquote=true,
  literate=
   {à}{{\`a}}1 {â}{{\^a}}1 {ä}{{\"a}}1
   {ç}{{\c c}}1
   {é}{{\'e}}1 {è}{{\`e}}1 {ê}{{\^e}}1 {ë}{{\"e}}1
   {î}{{\^\i}}1 {ï}{{\"\i}}1
   {ô}{{\^o}}1 {ö}{{\"o}}1
   {ù}{{\`u}}1 {û}{{\^u}}1 {ü}{{\"u}}1
   {À}{{\`A}}1 {Â}{{\^A}}1 {Ä}{{\"A}}1
   {Ç}{{\c C}}1
   {É}{{\'E}}1 {È}{{\`E}}1 {Ê}{{\^E}}1 {Ë}{{\"E}}1
   {Î}{{\^I}}1 {Ï}{{\"I}}1
   {Ô}{{\^O}}1 {Ö}{{\"O}}1
   {Ù}{{\`U}}1 {Û}{{\^U}}1 {Ü}{{\"U}}1
   {œ}{{\oe}}1 {Œ}{{\OE}}1
   {–}{{--}}1 {—}{{---}}1
}

\newcommand{\code}[1]{\texttt{#1}}
\newcommand{\pp}{\,points de pourcentage}
\newcommand{\taglabel}[1]{\scriptsize\sffamily\bfseries #1}

\newcounter{exercise}
\newcounter{extension}
\newcounter{logic}

\newcommand{\workbooktitle}[4]{%
  \setcounter{exercise}{0}%
  \setcounter{extension}{0}%
  \setcounter{logic}{0}%
  \fancyhead[L]{\sffamily\footnotesize #1 -- #2}%
  \begingroup\raggedright
  {\sffamily\Large\bfseries\color{ink}#1 -- #2\par}
  \vspace{0.25em}
  {\sffamily\small\color{muted}Python pour économistes -- Université de Lille, L2 Économie\par}
  \vspace{0.85em}
  {\sffamily\bfseries\color{accent}Question fil rouge.} #3\par
  \vspace{0.35em}
  {\sffamily\bfseries\color{green}Livrables.} #4\par
  \vspace{0.9em}
  {\color{line}\hrule height 0.45pt}
  \vspace{0.85em}
  \endgroup
}

\newenvironment{objectivebox}[1]{%
  \subsection*{#1}
  \vspace{-0.55em}
}{\vspace{0.15em}}

\newenvironment{routebox}[1]{%
  \subsection*{#1}
  \vspace{-0.45em}
}{\vspace{0.15em}}

\newenvironment{deliverablebox}[1]{%
  \subsection*{#1}
  \vspace{-0.45em}
}{\vspace{0.15em}}

\newenvironment{methodbox}[1]{%
  \subsection*{#1}
  \vspace{-0.45em}
}{\vspace{0.15em}}

\newenvironment{pathwaybox}[1]{%
  \par\addvspace{0.8em}
  {\sffamily\bfseries\color{blue}#1\par}
  \vspace{-0.25em}
  \begingroup\leftskip=0.85em\relax
  \noindent\color{ink}
}{\par\endgroup\addvspace{0.25em}}

\newenvironment{pseudocodeexercise}[1]{%
  \refstepcounter{logic}%
  \begin{tcolorbox}[
    enhanced,breakable,
    title={Logique \thelogic{} -- #1 \hfill\taglabel{PSEUDO-CODE}},
    colback=bluesoft,
    colframe=blue!55!black,
    coltitle=ink,
    colbacktitle=bluesoft,
    fonttitle=\sffamily\bfseries,
    boxrule=0.35pt,
    arc=0pt,
    left=10pt,right=10pt,top=7pt,bottom=7pt,
    before skip=8pt,after skip=8pt,
    borderline west={3pt}{0pt}{blue}
  ]%
}{\end{tcolorbox}}

\newenvironment{mandatoryexercise}[1]{%
  \refstepcounter{exercise}%
  \par\addvspace{1.05em}
  {\color{line}\hrule height 0.35pt}
  \vspace{0.55em}
  {\sffamily\large\bfseries\color{ink}Exercice \theexercise{} -- #1\par}
  {\taglabel{OBLIGATOIRE}\par}
  \vspace{0.25em}
}{\par\addvspace{0.65em}}

\newenvironment{extensionexercise}[1]{%
  \refstepcounter{extension}%
  \begin{tcolorbox}[
    enhanced,breakable,
    title={Pour aller plus loin \theextension{} -- #1},
    colback=ambersoft,
    colframe=amber!65!black,
    coltitle=ink,
    colbacktitle=ambersoft,
    fonttitle=\sffamily\bfseries,
    boxrule=0.35pt,
    arc=0pt,
    left=10pt,right=10pt,top=8pt,bottom=8pt,
    before skip=8pt,after skip=8pt,
    borderline west={3pt}{0pt}{amber}
  ]%
}{\end{tcolorbox}}

\newenvironment{checkpointbox}[1]{%
  \subsection*{#1}
  \vspace{-0.45em}
}{\vspace{0.15em}}

\newtcolorbox{takeawaybox}[1]{
  enhanced,breakable,
  title={#1},
  colback=greensoft,
  colframe=green,
  coltitle=ink,
  colbacktitle=greensoft,
  fonttitle=\sffamily\bfseries,
  boxrule=0.35pt,
  arc=0pt,
  left=10pt,right=10pt,top=8pt,bottom=8pt,
  before skip=10pt,after skip=8pt,
  borderline west={3pt}{0pt}{green}
}


\begin{document}

\workbooktitle
  {TD3}
  {Fonctions et comparaison de groupes}
  {Comment écrire des fonctions Python pour comparer des groupes de pays à partir d'indicateurs OWID ?}
  {Un fichier \code{td3\_nom\_prenom.py}, un tableau texte, un graphique \code{comparaison\_pib\_td3.png} et une interprétation de cinq lignes.}

\begin{objectivebox}{Objectifs du TD}
\begin{itemize}
  \item Distinguer paramètre, argument, affichage et valeur de retour.
  \item Écrire et tester des fonctions simples sur des listes numériques.
  \item Comparer deux groupes avec moyenne, médiane, minimum et maximum.
  \item Réutiliser une fonction pour produire un graphique \code{matplotlib}.
  \item Corriger des erreurs liées aux noms de variables et à \code{return}.
\end{itemize}
\end{objectivebox}

\begin{checkpointbox}{Progression du TD}
Les exercices 1 à 13 sont obligatoires. Les notions statistiques obligatoires sont : moyenne, médiane, minimum, maximum et proportion. Les quartiles et l'écart-type restent facultatifs. Le point Python central est la fonction : paramètre, appel, \code{return} et test.
\end{checkpointbox}

\begin{checkpointbox}{Source des données}
Les valeurs de PIB par habitant sont des extraits pédagogiques arrondis d'Our World in Data.
\begin{lstlisting}[language=terminal]
https://ourworldindata.org/grapher/gdp-per-capita-worldbank.csv
\end{lstlisting}
\end{checkpointbox}

\begin{pseudocodeexercise}{Automatiser une statistique}
\begin{lstlisting}
# Objectif : comparer deux listes de valeurs
# Entrees : liste_groupe_a, liste_groupe_b
# Definir une fonction moyenne(liste)
# Definir une fonction mediane(liste)
# Appliquer les fonctions aux deux groupes
# Comparer les resultats
# Sortie : tableau + phrase d'interpretation
\end{lstlisting}
\end{pseudocodeexercise}

\section*{A. Données et rappels}

\begin{mandatoryexercise}{Créer deux groupes de pays}
Créer deux listes de PIB par habitant, exprimées en milliers de dollars internationaux et arrondies pour l'exercice.
\begin{lstlisting}
groupe_europe = [46.0, 54.0, 34.0, 38.0, 22.0, 61.0]
groupe_monde = [46.0, 76.0, 19.0, 10.0, 7.0, 3.0]
\end{lstlisting}
Ajouter deux listes de noms pour documenter les valeurs :
\begin{lstlisting}
pays_europe = ["France", "Germany", "Spain", "Italy", "Poland", "Sweden"]
pays_monde = ["France", "United States", "China", "Brazil", "India", "Ethiopia"]
\end{lstlisting}
\end{mandatoryexercise}

\begin{mandatoryexercise}{Prédire quelques opérations sur les listes}
Sans écrire de fonction, prédire puis vérifier :
\begin{enumerate}
  \item \code{len(groupe\_europe)} ;
  \item \code{sum(groupe\_monde)} ;
  \item \code{sorted(groupe\_monde)} ;
  \item \code{groupe\_europe[0] > groupe\_monde[0]}.
\end{enumerate}
Ajouter une phrase qui explique pourquoi ces petits contrôles sont utiles avant de définir une fonction.
\end{mandatoryexercise}

\begin{mandatoryexercise}{Rappel intuitif}
Ajouter dans le script trois commentaires :
\begin{enumerate}
  \item la moyenne résume un niveau général ;
  \item la médiane est la valeur centrale après tri ;
  \item min et max donnent l'étendue des valeurs observées.
\end{enumerate}
\end{mandatoryexercise}

\section*{B. Fonctions}

\begin{mandatoryexercise}{Lire une fonction avant d'en écrire une}
Avant d'écrire des fonctions statistiques, lire ce petit exemple.
\begin{lstlisting}
def ajouter_trois(x):
    resultat = x + 3
    return resultat

print(ajouter_trois(10))
\end{lstlisting}
Répondre en commentaires :
\begin{enumerate}
  \item quel est le paramètre de la fonction ?
  \item quel est l'argument utilisé lors de l'appel ?
  \item quelle ligne renvoie une valeur ?
  \item quelle ligne affiche une valeur ?
\end{enumerate}
\end{mandatoryexercise}

\begin{mandatoryexercise}{Modifier une fonction courte}
Compléter puis modifier cette fonction.
\begin{lstlisting}
def convertir_en_dollars(valeur_milliers):
    return ...

print(convertir_en_dollars(46.0))
\end{lstlisting}
La fonction doit transformer \code{46.0} en \code{46000.0}. Modifier ensuite son nom ou son paramètre, puis vérifier que l'appel utilise exactement le même nom. L'objectif est de repérer les erreurs de nom avant les fonctions statistiques.
\end{mandatoryexercise}

\begin{mandatoryexercise}{Fonction moyenne}
Écrire une fonction \code{moyenne(valeurs)} qui renvoie la moyenne d'une liste. La tester sur les deux groupes.
\begin{lstlisting}
def moyenne(valeurs):
    return sum(valeurs) / len(valeurs)
\end{lstlisting}
La tester d'abord sur \code{[10, 12, 14]} pour vérifier que le résultat attendu est \code{12.0}. Afficher ensuite les résultats des deux groupes avec une unité : \code{milliers de dollars par habitant}.
\end{mandatoryexercise}

\begin{mandatoryexercise}{Fonction médiane guidée}
Pour une liste de longueur paire, on peut prendre la moyenne des deux valeurs centrales. Compléter la fonction :
\begin{lstlisting}
def mediane_paire(valeurs):
    valeurs_triees = sorted(valeurs)
    n = len(valeurs_triees)
    milieu_droit = ...
    milieu_gauche = ...
    return ...
\end{lstlisting}
Tester d'abord la fonction sur \code{[8, 10, 12, 14]} : la médiane attendue est \code{11.0}. Tester ensuite la fonction sur les deux groupes.
\end{mandatoryexercise}

\begin{mandatoryexercise}{Débogage de fonction}
Les deux fonctions suivantes ressemblent à des solutions, mais elles posent problème. Identifier l'erreur, formuler une hypothèse, puis écrire la correction.
\begin{lstlisting}
def moyenne_bug(valeurs):
    resultat = sum(valeur) / len(valeurs)
    return resultat

def moyenne_sans_return(valeurs):
    print(sum(valeurs) / len(valeurs))

m = moyenne_sans_return([10, 12, 14])
print(m + 1)
\end{lstlisting}
Indice : le premier cas concerne un nom de variable ; le second permet de distinguer afficher une valeur et renvoyer une valeur.
\end{mandatoryexercise}

\begin{mandatoryexercise}{Min, max et proportion}
Pour chaque groupe, calculer :
\begin{enumerate}
  \item la valeur minimale ;
  \item la valeur maximale ;
  \item la proportion de pays au-dessus de 30 milliers de dollars par habitant.
\end{enumerate}
Conseil : utiliser un compteur et une boucle.
\end{mandatoryexercise}

\section*{C. Comparer et interpréter}

\begin{pseudocodeexercise}{Comparer deux groupes par un graphique}
\begin{lstlisting}
# Objectif : comparer visuellement deux resultats
# Entree : deux listes de valeurs et une fonction moyenne
# 1. Calculer la moyenne du groupe Europe
# 2. Calculer la moyenne du groupe Monde
# 3. Placer les deux resultats dans une liste
# 4. Tracer une barre par groupe
# 5. Sauvegarder la figure
# Sortie : comparaison lisible
\end{lstlisting}
\end{pseudocodeexercise}

\begin{mandatoryexercise}{Graphique de comparaison}
Utiliser les fonctions écrites plus haut pour construire un graphique simple.
\begin{lstlisting}
import matplotlib.pyplot as plt

groupes = ["Europe", "Monde"]
moyennes = [moyenne(groupe_europe), moyenne(groupe_monde)]

plt.figure()
plt.bar(groupes, moyennes)
plt.ylabel("Milliers de dollars par habitant")
plt.title("PIB par habitant moyen - extrait OWID")
plt.tight_layout()
plt.savefig("comparaison_pib_td3.png", dpi=300)
\end{lstlisting}
Modifier ensuite le code pour tracer les médianes à la place des moyennes.
\end{mandatoryexercise}

\begin{mandatoryexercise}{Déboguer un graphique}
Voici deux fragments qui ressemblent au code précédent mais posent problème. Identifier l'erreur probable, puis proposer une correction.
\begin{lstlisting}
# Fragment A
groupes = ["Europe", "Monde"]
moyennes = [moyenne(groupe_europe)]
plt.bar(groupes, moyennes)

# Fragment B
plt.bar(groupes, medianes)
\end{lstlisting}
Dans le fragment A, comparer la longueur des deux listes. Dans le fragment B, vérifier si la variable \code{medianes} a bien été créée avant son utilisation.
\end{mandatoryexercise}

\begin{mandatoryexercise}{Tableau texte}
Afficher un petit tableau texte.
Il doit contenir les colonnes \code{groupe}, \code{moyenne}, \code{mediane}, \code{min}, \code{max} et \code{part\_au\_dessus\_30}.
L'objectif est la lisibilité, pas la perfection typographique.
\end{mandatoryexercise}

\begin{mandatoryexercise}{Question économique}
Ajouter cinq commentaires à la fin du script :
\begin{enumerate}
  \item quel groupe a la moyenne la plus élevée ?
  \item quel groupe a la médiane la plus élevée ?
  \item pourquoi moyenne et médiane peuvent-elles raconter deux choses différentes ?
  \item pourquoi le choix des pays influence-t-il fortement le résultat ?
  \item quelle limite voyez-vous à travailler sur seulement six pays par groupe ?
\end{enumerate}
\end{mandatoryexercise}

\section*{D. Entraînement facultatif}

\begin{extensionexercise}{Statistiques supplémentaires sans pression}
Choisir au moins deux questions si le parcours obligatoire est terminé.
\begin{enumerate}
  \item Ajouter un pays dans chaque groupe et vérifier que les fonctions marchent encore.
  \item Calculer l'étendue : \code{max - min}.
  \item Écrire une fonction \code{proportion\_au\_dessus(valeurs, seuil)}.
  \item Tester le seuil de 50 au lieu de 30.
  \item Écrire une fonction \code{resume(valeurs)} qui affiche moyenne, médiane, min et max.
  \item Comparer la moyenne avant et après ajout d'une valeur très élevée.
  \item Calculer un premier quartile avec une méthode simple de votre choix.
  \item Calculer l'écart-type à partir d'une formule fournie par l'enseignant.
  \item Écrire une phrase expliquant pourquoi l'écart-type mesure la dispersion.
  \item Transformer les valeurs en dollars au lieu de milliers de dollars.
\end{enumerate}
\end{extensionexercise}

\begin{takeawaybox}{À retenir}
\begin{itemize}
  \item Une fonction sert à nommer une suite d'opérations que l'on veut réutiliser.
  \item Le paramètre est écrit dans la définition ; l'argument est la valeur fournie lors de l'appel.
  \item \code{return} renvoie une valeur réutilisable, alors que \code{print()} affiche seulement.
  \item Une fonction testée peut servir ensuite dans un tableau ou un graphique.
  \item Tester une fonction sur un cas simple permet de détecter une erreur avant de l'appliquer aux données.
\end{itemize}
\end{takeawaybox}

\end{document}
